\section{Real-System Validation}
\label{sec:real-system}

\subsection{Objective}
\label{sec:real-system-objective}

This analysis addresses RQ6 (\S\ref{sec:rq}). We do not claim that simulated absolute latency or throughput values match real-hardware values; simulator latency is not assumed to transfer directly to real-engine latency, nor is the simulator's own load-calibration reference $\lambda_{\mathrm{ref}}$ (\S\ref{sec:calibration}) expected to transfer to a physical engine's saturation point. The sole objective of this validation is to test agreement in \emph{relative policy ordering} and \emph{reversal direction} between the simulator and a real vLLM deployment on a targeted set of validation cases.

\subsection{Experimental Design}
\label{sec:real-system-design}

The real-system validation follows a preregistered, window-resampled evaluation design structured to mirror the simulated campaign without requiring exhaustive duplication. The evaluation protocol is defined as follows:

\begin{itemize}
  \item \textbf{Inferential Unit:} The statistical unit of analysis is the workload window, preserving the methodology established in \S\ref{sec:methodology}.
  \item \textbf{Scale and Execution:} The design evaluates $40$ independent workload windows per source, each containing $200$ requests. Exactly one scientific execution is performed per (policy, source, window) cell, yielding $240$ total scientific tasks ($2$ policies $\times$ $3$ sources $\times$ $40$ windows).
  \item \textbf{Uncertainty Quantification:} Uncertainty is quantified via window-level bootstrap resampling ($\ge 2{,}000$ resamples) over the $40$ independent windows per source.
  \item \textbf{Operating Regions:} Workloads are calibrated against the real engine's observed saturation point to define load regions independently of the simulator's calibration.
  \item \textbf{Software and Hardware:} Every execution serves \texttt{Qwen/Qwen2.5-0.5B-Instruct} on \texttt{vllm==0.27.1} (\texttt{torch==2.13.0}), one instance per task on a single NVIDIA A100-SXM4-80GB GPU (Slurm resource label \texttt{gres=gpu:a100\_10g}, unchanged from the validated load-calibration campaign). This label is a Slurm resource tag, not evidence of a fractional or MIG GPU partition; the allocated device is a full 80\,GB A100. Tokenizer revision, model dtype, tensor-parallel degree, and maximum model length were not separately recorded in the frozen execution manifest; rather than invent values, we do not report them here.
  \item \textbf{Ordering and Provenance:} Each (policy, source, window) cell is executed exactly once; with one execution per cell there is no repeated-measures ordering to alternate. Cells are assigned deterministically to Slurm array indices by a fixed (source, window, policy) sort, with execution concurrency limited to four simultaneous array tasks. Every execution records its exact git commit, hardware environment (per-task GPU query), software versions, and workload/calibration-manifest hash identity to ensure reproducibility.
\end{itemize}

\subsection{Case Selection}
\label{sec:real-system-case-selection}

Validation cases were selected mechanically from the validated results presented in \S\ref{sec:reversals}, applying a preregistered rule that prefers the pairwise reversal with the largest effect size and a representative stable-ordering control.

\begin{itemize}
  \item \textbf{Reversal Case:} \texttt{slai\_faithful} versus \texttt{vllm\_faithful}, comparing the Azure and BurstGPT sources at high load. This case exhibited the largest-effect-size statistically supported practical reversal in simulation (a simulator difference of $+0.520$ for Azure favoring \texttt{slai\_faithful}, and $-0.348$ for BurstGPT favoring \texttt{vllm\_faithful}).
  \item \textbf{Stable Control:} Azure versus Bailian/Qwen at high load, chosen for exhibiting perfect rank agreement in simulation.
\end{itemize}

Evaluating this selection on physical hardware requires mapping the simulated mechanisms to their real-engine equivalents. While \texttt{vllm\_faithful} maps directly to native vLLM execution paths, the SLO-aware SLAI/RAD mechanism requires a custom vLLM scheduling plugin to realize its semantics on real hardware.

\subsection{Results}
\label{sec:real-system-result}

The physical-hardware load-calibration prerequisite for these validation cases has been executed (120/120 terminal outcomes: 43 converged, 77 lower-bound-already-violating), establishing the necessary injection rates for the deployment environment. All $240$ scientific evaluation tasks ($2$ policies $\times$ $3$ sources $\times$ $40$ windows) subsequently completed on the software/hardware environment above (\S\ref{sec:real-system-design}), with zero failed, timed-out, or partial cells.

Table~\ref{tab:rq6-results} reports the \texttt{slai\_faithful}-minus-\texttt{vllm\_faithful} arrival-normalized-weighted-goodput (ANWG) effect for each source, using the same window-level paired bootstrap as \S\ref{sec:uncertainty} ($2{,}000$ resamples over the $40$ independent windows/source, 95\% CI). All three effects are statistically supported (CI excludes zero) and favor \texttt{vllm\_faithful}.

\begin{table}[t]
\centering
\caption{Real-vLLM RQ6 outcomes against the simulator-selected cases (\texttt{HIGH\_PRESSURE}, 40 paired windows/source, 95\% window-level bootstrap CI, $2{,}000$ resamples). Effect is \texttt{slai\_faithful}$-$\texttt{vllm\_faithful} ANWG.}
\label{tab:rq6-results}
\small
\setlength{\tabcolsep}{4.5pt}
\begin{tabular}{lccccl}
\toprule
Source & Sim.\ winner & Effect & 95\% CI & Real winner & Agreement \\
\midrule
Azure (rev.) & SLAI ($+0.520$) & $-0.024$ & $[-0.027,-0.020]$ & vLLM & Disagrees \\
BurstGPT (rev.) & vLLM ($-0.348$) & $-0.200$ & $[-0.264,-0.144]$ & vLLM & Agrees \\
Bailian (ctrl.) & n/a$^{*}$ & $-0.026$ & $[-0.030,-0.022]$ & vLLM & Consistent \\
\bottomrule
\end{tabular}
\vspace{2pt}

{\footnotesize $^{*}$The simulator's stable-control quantity ($\tau_b = 1.00$) is an
11-policy full-ranking correlation, not a two-policy \texttt{slai\_faithful}-vs-\texttt{vllm\_faithful}
winner; no directly comparable simulator winner is defined for this row. The
relevant real-system check is instead whether the two-policy ordering agrees
between the Azure and Bailian/Qwen conditions, which it does (both favor
\texttt{vllm\_faithful}).}
\end{table}

The simulator-predicted Azure/BurstGPT reversal did not reproduce: \texttt{vllm\_faithful} won both conditions on real hardware, so the predicted sign flip between Azure and BurstGPT was not observed. BurstGPT's real winner matches the simulator's selection; Azure's does not -- the physical Azure comparison changed from a simulated \texttt{slai\_faithful} advantage to a small but statistically supported \texttt{vllm\_faithful} advantage. In contrast, the Azure/Bailian-Qwen stable-control comparison retained a consistent \texttt{vllm\_faithful}-favoring ordering across both real-system conditions (effect magnitudes $-0.0235$ and $-0.0256$), reproducing the qualitative stability that comparison was selected to test, independent of which specific policy wins.

Statistical support is not the same as practical magnitude. Read against
this study's 10\% practical-margin convention (\S\ref{sec:uncertainty}),
the three real-system effects vary substantially once expressed relative
to the losing policy's own mean ANWG (frozen descriptive aggregates,
\texttt{docs/RQ6\_PUBLIC\_RESULT\_PROVENANCE.md}): Azure's margin is
$(0.7700-0.7465)/0.7465 \approx 3.1\%$ and Bailian/Qwen's is
$(0.84525-0.819625)/0.819625 \approx 3.1\%$ -- both statistically
supported but below the 10\% convention used elsewhere in this study to
separate practically meaningful from microscopic effects. BurstGPT's
margin, $(0.316625-0.116375)/0.116375 \approx 172\%$, substantially
exceeds it. This is a descriptive comparison against an existing
convention, not a new inferential analysis: \S\ref{sec:uncertainty}'s
five-class taxonomy is defined for a pairwise comparison across two
conditions of the same source and is not applied here to a single
physical-source row.

This selected-case result exposes a simulator-fidelity boundary rather
than establishing a general finding: the Azure/BurstGPT reversal that the
simulator identified as this study's largest-effect-size case did not
transfer to physical execution, while the paired stable-control comparison
retained its qualitative ordering (\S\ref{sec:threats} scopes what this
design does and does not establish; \S\ref{sec:discussion} interprets the
two outcomes together).
